\documentclass[10pt]{article}
\providecommand{\FIGDIR}{fig}
\newcommand{\DIGESTDATE}{28 July 2026}
% ===================== PM Research Watch - shared preamble =====================
\usepackage[T1]{fontenc}
\usepackage[utf8]{inputenc}
\usepackage{amsmath,amssymb}
\usepackage{tgpagella}
\usepackage{tgheros}
\renewcommand{\sfdefault}{qhv}
\renewcommand{\ttdefault}{lmtt}
\usepackage{microtype}
\usepackage[english]{babel}

\usepackage[a4paper,top=17mm,bottom=17mm,left=17mm,right=17mm,
            headsep=6mm,footskip=10mm]{geometry}
\usepackage{graphicx}
\usepackage{xcolor}
\usepackage{colortbl}
\usepackage{tikz}
\usetikzlibrary{calc}
\usepackage[most]{tcolorbox}
\usepackage{titlesec}
\usepackage{fancyhdr}
\usepackage{booktabs}
\usepackage{array}
\usepackage{longtable}
\usepackage{enumitem}
\usepackage{multicol}
\usepackage{ragged2e}
\usepackage[hidelinks]{hyperref}

% ------------------------------------------------------------------ palette
\definecolor{Ink}    {HTML}{14262E}
\definecolor{Deep}   {HTML}{0F4C5C}
\definecolor{Teal}   {HTML}{1F7A8C}
\definecolor{Sky}    {HTML}{4FA3C4}
\definecolor{Sage}   {HTML}{6FA287}
\definecolor{Amber}  {HTML}{E3A008}
\definecolor{Clay}   {HTML}{B5643C}
\definecolor{Coral}  {HTML}{D1495B}
\definecolor{Violet} {HTML}{6A5B9E}
\definecolor{Slate}  {HTML}{4A6572}
\definecolor{Paper}  {HTML}{FBF9F5}
\definecolor{Mist}   {HTML}{EEEAE2}
\definecolor{Rule}   {HTML}{D9D3C9}

\pagecolor{Paper}
\color{Ink}

\hypersetup{colorlinks=true, linkcolor=Deep, urlcolor=Teal, citecolor=Teal,
            pdftitle={PM Research Watch}, pdfauthor={Automated daily digest}}

% ------------------------------------------------------------------ headings
\titleformat{\section}
  {\sffamily\bfseries\large\color{Deep}}{}{0pt}
  {}[{\vspace{1.2mm}\color{Amber}\titlerule[1.1pt]}]
\titlespacing*{\section}{0pt}{5.5mm}{2.8mm}

\titleformat{\subsection}
  {\sffamily\bfseries\normalsize\color{Teal}}{}{0pt}{}
\titlespacing*{\subsection}{0pt}{3.6mm}{1.4mm}

\setlist[itemize]{leftmargin=4.2mm, itemsep=1.1mm, topsep=1.3mm, parsep=0pt,
                  label=\textcolor{Teal}{\footnotesize$\blacktriangleright$}}

% ------------------------------------------------------------------ header/footer
\pagestyle{fancy}
\fancyhf{}
\renewcommand{\headrulewidth}{0pt}
\renewcommand{\footrulewidth}{0pt}
\fancyhead[L]{\sffamily\scriptsize\color{Slate}PM RESEARCH WATCH\;\textcolor{Rule}{\textbar}\;\DIGESTDATE}
\fancyhead[R]{\sffamily\scriptsize\color{Slate}Particulate matter monitoring \& health effects}
\fancyfoot[C]{\sffamily\scriptsize\color{Slate}\thepage}

% ------------------------------------------------------------------ components
% metric tile
\newcommand{\metric}[3]{%
  \begin{tcolorbox}[enhanced, colback=white, colframe=Rule, boxrule=0.5pt,
    arc=1.6mm, left=2.2mm, right=2.2mm, top=1.8mm, bottom=1.8mm,
    borderline west={1.7mm}{0pt}{#3}, nobeforeafter, valign=center,
    height=17mm, width=\linewidth]
    \RaggedRight
    {\sffamily\bfseries\fontsize{17}{17}\selectfont\color{#3}#1}\\[0.5mm]
    {\sffamily\fontsize{6.5}{7.4}\selectfont\color{Slate}#2}
  \end{tcolorbox}}

% pull-quote / key insight box
\newtcolorbox{keybox}[1][]{enhanced, colback=Mist, colframe=Mist,
  arc=1.6mm, left=3.2mm, right=3.2mm, top=2.4mm, bottom=2.4mm,
  borderline west={1.7mm}{0pt}{Amber}, #1}

% subtopic band
\newcommand{\band}[2]{%
  \vspace{3.2mm}\noindent
  \begin{tcolorbox}[enhanced, colback=#2, colframe=#2, arc=1.2mm,
    left=2.6mm, right=2.6mm, top=1.1mm, bottom=1.1mm, boxrule=0pt,
    width=\linewidth, nobeforeafter]
    {\sffamily\bfseries\small\color{white}#1}
  \end{tcolorbox}\vspace{1.6mm}}

% one paper entry: {short cite}{title}{journal}{doi}{body}
\newcommand{\paperentry}[5]{%
  \par\vspace{1.9mm}\noindent
  {\sffamily\bfseries\footnotesize\color{Ink}#1}\;%
  {\sffamily\fontsize{7}{8}\selectfont\color{Slate}#3}\\[0.4mm]
  {\footnotesize\itshape\color{Deep}#2}\\[0.6mm]
  {\footnotesize #4}\\[0.3mm]
  {\fontsize{6.6}{7.6}\selectfont\color{Slate}\texttt{doi:}\,\href{https://doi.org/#5}{\texttt{\detokenize{#5}}}}
  \par}

% caption style
\newcommand{\figcap}[1]{\par\vspace{1.0mm}
  {\fontsize{7.2}{8.6}\selectfont\color{Slate}\textsf{\textbf{Fig.}}\; #1\par}\vspace{1.5mm}}

\providecommand{\permil}{\textperthousand}

\setlength{\parindent}{0pt}
\setlength{\parskip}{0pt}
\renewcommand{\arraystretch}{1.18}


\begin{document}

% ==================================================================== MASTHEAD
\thispagestyle{empty}
\begin{tikzpicture}[remember picture, overlay]
  \fill[Deep] (current page.north west) rectangle
      ($(current page.north east)+(0,-67mm)$);
  \fill[Amber] ($(current page.north west)+(0,-67mm)$) rectangle
      ($(current page.north east)+(0,-68.4mm)$);
  \foreach \x/\y/\r in {22/12/0.9, 41/26/1.5, 63/9/0.7, 88/31/1.1, 112/17/1.9,
                        138/29/0.8, 159/13/1.3, 178/54/1.0, 196/21/1.6,
                        30/44/0.6, 74/57/1.2, 125/48/0.9, 168/6/0.8, 191/59/0.7,
                        52/18/1.0, 100/60/0.6, 148/20/1.1, 66/49/1.4, 205/45/0.8}
    \fill[Sky, opacity=0.30] ($(current page.north west)+(\x mm,-\y mm)$) circle (\r mm);
\end{tikzpicture}

\vspace*{4mm}
{\sffamily\fontsize{7.6}{9}\selectfont\color{Sky}%
 AUTOMATED DAILY LITERATURE DIGEST\;\textbullet\;PhD RESEARCH WATCH\par}
\vspace{2.2mm}
{\sffamily\bfseries\fontsize{27}{29}\selectfont\color{white}%
 Particulate Matter\par}
{\sffamily\fontsize{27}{29}\selectfont\color{Sky}%
 Monitoring \& Health Effects\par}
\vspace{5.0mm}
{\sffamily\fontsize{8.4}{10}\selectfont\color{white}%
 Issue 2026-07-28\;\;\textcolor{Amber}{\textbar}\;\;Entry window 28 July 2026 (PubMed EDAT);
 Europe PMC creation 26--28 July\;\;%
 \textcolor{Amber}{\textbar}\;\;21 screened, 19 in scope\par}
\vspace{18mm}

% -------------------------------------------------------------- metric strip
\noindent
\begin{minipage}[t]{0.190\linewidth}\metric{19}{RECORDS IN SCOPE\\(2 EXCLUDED)}{Deep}\end{minipage}\hfill
\begin{minipage}[t]{0.190\linewidth}\metric{10}{SUBTOPIC CLUSTERS\\ALL REPRESENTED}{Teal}\end{minipage}\hfill
\begin{minipage}[t]{0.190\linewidth}\metric{8}{RECORDS FROM CHINA\\(42\% OF CORPUS)}{Sky}\end{minipage}\hfill
\begin{minipage}[t]{0.190\linewidth}\metric{2}{ESTIMATES REPORTED\\WITH A 95\% CI}{Amber}\end{minipage}\hfill
\begin{minipage}[t]{0.190\linewidth}\metric{12}{RECORDS IN MDPI\\TITLES (63\%)}{Clay}\end{minipage}

\vspace{6mm}

% -------------------------------------------------------------- signal of day
\section*{\sffamily\bfseries\color{Deep}Signal of the day}
\vspace{-1.5mm}{\color{Amber}\rule{\linewidth}{1.1pt}}\vspace{2.5mm}

\begin{keybox}
{\sffamily\bfseries\small\color{Deep}The binding constraint on African exposure
epidemiology is record linkage, not instrumentation}\\[1.6mm]
{\footnotesize
The AIRCLARE-Ghana pilot (Research Square preprint, \texttt{rs-10441248}) is the first
time-stratified case-crossover study of short-term particulate exposure in Ghana. Two
years of \emph{federal-equivalent-method} PM$_{2.5}$/PM$_{10}$ monitoring were linked to
213{,}241 outpatient records at the University of Ghana Hospital, yielding 9{,}511 acute
respiratory presentations. Each IQR rise in PM$_{2.5}$ (14.70\,$\mu$g/m$^3$) raised the
odds of an upper-respiratory presentation by
\textbf{3.8\% (OR 1.038, 95\% CI 1.010--1.068)}; PM$_{10}$ (IQR 51.49\,$\mu$g/m$^3$) gave
\textbf{OR 1.030 (1.008--1.052)}. The lag structure peaks at day 1
(OR $\approx$ 1.07 / 1.06) and persists to a cumulative 0--7\,d
(OR $\approx$ 1.09 / 1.06), with the strongest signal during the dry Harmattan season and
among women.\\[1.4mm]
Two things matter here for a distributed-sensing programme. First, the reported
\textbf{96.03\% monitoring-stream completeness after QC} is the study's real contribution:
it shows the reference-grade data layer is not the failure point in a low-resource
deployment --- the pipeline that joins it to clinical records is. Second, the effect sizes
are small enough that they are only recoverable because exposure misclassification was
minimised by a single, co-located, high-completeness monitor serving a geographically
compact catchment. Any attempt to scale this design across Accra with a distributed
low-cost network will face the opposite trade-off: better spatial resolution, worse and
drift-dependent measurement error. Pneumonia and asthma associations were directionally
consistent but null, which is what an underpowered pilot at these effect sizes should
produce, and the authors say so.\par}
\end{keybox}

\vspace{2mm}
{\fontsize{7.6}{9.2}\selectfont\color{Slate}Caveats stated plainly: this is a
\textbf{preprint}, not peer reviewed; the catchment is a university hospital and student
clinic, so the population skews young, educated and urban and the 925.7-per-1{,}000 URI
rate is not a community incidence; a single monitoring site is treated as the exposure for
all subjects, so within-city gradients are absorbed into error; and the six-stage ICD-10
filtering cascade is a bespoke, unvalidated case definition whose sensitivity/specificity
are unreported. The case-crossover design does control time-invariant confounding by
construction.\par}

\vspace{4mm}

% -------------------------------------------------------------- what changed
\section*{\sffamily\bfseries\color{Deep}What changed today}
\vspace{-1.5mm}{\color{Amber}\rule{\linewidth}{1.1pt}}\vspace{2.5mm}

\begin{itemize}
\item \textbf{Sub-Saharan Africa enters the corpus.} Yesterday's 33-record issue contained
zero records from the region; today two of nineteen (Ghana case-crossover; Uganda scoping
review) originate there. Both are explicitly about monitoring infrastructure rather than
biology, which is the correct order of operations for a region with essentially no
long-term reference network.

\item \textbf{Quantitative reporting quality collapsed.} Yesterday twelve effect estimates
with intervals were extractable from 33 records. Today \textbf{two of nineteen} records
report a point estimate with a 95\% CI. Several report bare point estimates
(PD-L1 OR $=$ 1.035, $p=0.017$), ranges across unspecified model sets
(sleep disorders, ``OR 1.08 to 3.20''), or standardised regression coefficients with
significance stars only (spatial Durbin, 0.028\textsuperscript{*}). This is not a
sampling artefact --- it tracks the journal-composition shift below.

\item \textbf{Journal composition shifted sharply.} Twelve of nineteen records (63\%) are
in MDPI titles (\textit{Toxics} alone contributes seven). Reference-grade cohort work in
\textit{JAMA Netw Open}, \textit{J Am Heart Assoc} and \textit{Diabetologia}, which
carried yesterday's issue, is absent. Weight the corpus accordingly: today's tier-A count
is three, against nine yesterday.

\item \textbf{Instrumentation produced nothing.} Zero records in this entry window concern
low-cost sensor characterisation, co-location calibration, drift correction, network
siting, or satellite/model fusion. The \textit{Sensing} cluster is filled by a bioaerosol
surveillance campaign and a bibliometric review. A supplementary sweep found relevant 2026
work (e.g.\ AMT 19, 1293) but it falls outside the entry window and belongs to no issue of
this digest. The gap is real, not a query failure: PubMed does not index
\textit{Atmos.\ Meas.\ Tech.}\ or \textit{Atmos.\ Chem.\ Phys.}, Europe PMC returned no
creations for 28 July, and OpenAlex has moved its \texttt{created\_date} filter behind a
paid plan --- see \emph{Provenance}.

\item \textbf{Microenvironmental measurement is the one growth area.} Five of nineteen are
measurement campaigns and all five are indoor or near-source: subway platforms under
different screen-door geometries, charcoal-barbecue dining rooms, an urban supersite, and
SARS-CoV-2 RNA recovered from ambient PM$_{2.5}$ filters. The common structure is a fixed
sampler in a small, high-concentration volume --- exactly the regime where a calibrated
low-cost node would be adequate and a reference sampler is overkill.
\end{itemize}

\clearpage

% ==================================================================== ANALYTICS
\section*{\sffamily\bfseries\color{Deep}Corpus analytics}
\vspace{-1.5mm}{\color{Amber}\rule{\linewidth}{1.1pt}}\vspace{3mm}

\noindent\includegraphics[width=\linewidth]{\FIGDIR/f1_subtopics.png}
\figcap{Subtopic assignment of the 19 in-scope records; single-label by dominant
contribution. Unlike yesterday's issue there is no dominant cluster: \textit{Burden,
policy \& mitigation} and \textit{Other clinical endpoints} lead with three each and six
clusters carry one or two. A flat distribution across a small $n$ carries little
information about the field --- it mostly reflects that one publisher's mid-month batch
supplied most of the corpus.}

\vspace{2mm}
\noindent\includegraphics[width=\linewidth]{\FIGDIR/f2_design_endpoint.png}
\figcap{Left --- study architecture. Measurement campaigns are the largest single class
(5/19, 26\%), a reversal of the review-heavy 27 July issue. Only one prospective cohort
appears, and no randomised or controlled-exposure human study. Right --- endpoint by organ
system. Oncologic, carcinogenic-risk and respiratory endpoints tie at two; the long tail
of single-record endpoints (sleep, auditory, placental, microenvironment, aerosol
chemistry) is characteristic of a corpus assembled from a general-toxicology journal
issue rather than from a coherent research front.}

\vspace{2mm}
\noindent\includegraphics[width=\linewidth]{\FIGDIR/f3_metric_geography.png}
\figcap{Left --- particle metric. PM$_{2.5}$ alone is the exposure in nine records and
appears jointly with PM$_{10}$ in five more. Ultrafine particles are again the explicit
subject of exactly one record, and that record measures no particles at all (it is a
spatial-econometric inference about UFP sources from industrial-output statistics). Right
--- geographic provenance. China supplies eight (42\%), Southeast Asia three (all Thailand,
all Chiang Mai haze-season work), and sub-Saharan Africa two. No record from South Asia.}

\vspace{2mm}
\noindent\includegraphics[width=\linewidth]{\FIGDIR/f6_lifecourse.png}
\figcap{Life-course windows addressed; counts are non-exclusive and reflect the
population actually studied, not the population the authors generalise to. The mass sits
in working-age adulthood (8) because today's measurement campaigns sample occupational and
commuter microenvironments. The in-utero window, dominant yesterday at five records, drops
to one --- and that one is a cell-line experiment in which the PM arm was null.}

\vspace{2mm}
\noindent\includegraphics[width=\linewidth]{\FIGDIR/f4_heatmap.png}
\figcap{Subtopic $\times$ study-architecture cross-tabulation. The diagnostic feature is
how empty it is: 19 records occupy 14 of 70 cells, with no cell above two. Note the
\textit{Sensing} row --- one measurement campaign, one review, and no modelling and no
linkage to any health endpoint. The \textit{Occupational \& indoor} row is purely
measurement, with no exposure--response analysis attached to either campaign. The
instrumentation/epidemiology silo flagged on 27 July persists unchanged.}

\vspace{2mm}
\noindent\includegraphics[width=\linewidth]{\FIGDIR/f5_forest.png}
\figcap{Every quantitative association extractable from today's corpus, on a common
logarithmic axis. \textbf{Only the two bottom-most estimates carry a reported interval};
markers drawn without whiskers are point estimates for which the source published no CI,
and they should be read as unfalsifiable until the full text is available. Note also the
axis range: all estimates lie between 1.03 and 1.09, an order of magnitude tighter than
yesterday's spread (1.03--3.16), because today's corpus contains no high-contrast
exposure comparison and no rare-outcome registry study. The
``sleep disorder'' marker is the \emph{lower bound} of a reported OR range (1.08--3.20)
whose model set is unspecified in the abstract; it is plotted at that bound deliberately.}

\clearpage

% ==================================================================== DIGEST
\section*{\sffamily\bfseries\color{Deep}Paper-level digest}
\vspace{-1.5mm}{\color{Amber}\rule{\linewidth}{1.1pt}}\vspace{1mm}

{\fontsize{7.4}{9}\selectfont\color{Slate}Ordered by cluster. Each entry gives the
methodological core, the finding that matters, the limitation that actually threatens the
inference, and the read-across to distributed low-cost PM sensing.\par}

\band{Neuro, cognition \& mental health\hfill 1 record}{Deep}

\paperentry{Lu et al.\ 2026b}
{Subchronic PM$_{2.5}$ exposure exacerbates chronic stress-induced depression-like pathology in mice}
{Environ Res}
{Factorial murine design crossing subchronic PM$_{2.5}$ inhalation with chronic
unpredictable mild stress (CUMS), read out by behavioural batteries, hippocampal
histopathology, transcriptomics, electrophysiology and synaptic-plasticity markers.
The result that matters is the \textbf{interaction, not the main effect}: PM$_{2.5}$
alone produced no depression- or anxiety-like behaviour, only modest neurobiological
change, whereas PM$_{2.5}$\,+\,CUMS worsened anhedonia, behavioural despair and anxiety
relative to CUMS alone, with greater hippocampal neuroinflammation, redox disruption and
impaired synaptic plasticity. This supports an effect-modifier rather than a direct-cause
model. The abstract reports no exposure concentration, no particle characterisation and no
effect sizes, so dose--response is unassessable and the inhalation regime cannot be
compared with ambient exposure; ``real-ambient'' status is not claimed. For sensing, the
implication is that a monitoring design aimed at mental-health endpoints needs to resolve
\emph{co-occurrence} of exposure and psychosocial stressor, which is a personal-monitoring
problem, not a fixed-site one.}
{10.1016/j.envres.2026.125327}

\band{Cardiovascular \& metabolic\hfill 1 record}{Teal}

\paperentry{Zhou et al.\ 2026}
{Air pollution and hypertension: an umbrella review of systematic reviews and meta-analyses}
{iScience}
{Umbrella review over three databases to September 2025, AMSTAR-2-appraised, retaining
eleven meta-analyses (eight on essential hypertension, three on gestational hypertension).
Conclusion: cumulative evidence supports PM$_{2.5}$ and PM$_{10}$ raising essential
hypertension risk, while the gestational-hypertension association remains inconclusive on
three studies. The value of this record is as a citation anchor and a map of
between-review heterogeneity; the abstract reports \textbf{no pooled estimate, no
$I^2$, and no AMSTAR-2 distribution}, so the strength of the ``plausible role'' claim
cannot be checked without the full text. Umbrella reviews also inherit the primary
literature's exposure-model heterogeneity --- land-use regression, satellite-derived and
nearest-monitor assignments are pooled as though interchangeable, which is precisely the
assumption a co-located sensor network is positioned to test.}
{10.1016/j.isci.2026.116811}

\band{Reproductive \& developmental\hfill 1 record}{Sky}

\paperentry{Liu et al.\ 2026}
{Deoxynivalenol, cytokines and particulate matter in human trophoblast cell models}
{Int J Mol Sci}
{Head-to-head sensitivity comparison of HTR-8/SVneo and BeWo trophoblast lines against
three maternal-exposome stimuli --- the mycotoxin deoxynivalenol, a TNF-$\alpha$/IL-1$\beta$/
IFN-$\gamma$ cocktail, and PM$_{2.5}$ --- at two penicillin/streptomycin concentrations,
24\,h exposure. \textbf{The PM arm was null}: no change in IL-6, IL-8 or hCG in either
line, while DON and the cytokine cocktail moved IL-6/IL-8, hCG and junctional markers
(ZO-1, OCLD, CLDN-3/4, E-CAD). A negative result in a 24\,h submerged monoculture at an
unreported dose is weak evidence of placental safety and strong evidence that these two
lines are the wrong model for particle toxicity --- there is no air--liquid interface, no
deposition physics and no particle characterisation. Worth logging precisely because
prenatal-PM epidemiology keeps citing trophoblast inflammation as its mechanism.}
{10.3390/ijms27146434}

\band{Respiratory \& allergic\hfill 2 records}{Sage}

\paperentry{AIRCLARE-Ghana 2026 (preprint)}
{Short-term ambient PM$_{2.5}$/PM$_{10}$ and acute respiratory morbidity in Ghana}
{Research Square}
{See \textit{Signal of the day}. Time-stratified case-crossover, Jan 2023--Dec 2024,
9{,}511 acute respiratory presentations from 213{,}241 outpatient records, conditional
logistic regression adjusted for temperature and relative humidity.
\textbf{URI: OR 1.038 (1.010--1.068) per IQR PM$_{2.5}$; OR 1.030 (1.008--1.052) per IQR
PM$_{10}$}, peaking at lag 1 and persisting to cumulative lag 0--7\,d. Pneumonia null,
asthma null. Threat to inference: one monitoring site stands in for a whole catchment, so
Berkson-type and classical exposure error coexist and the true effect is plausibly larger;
the case definition is bespoke and unvalidated; and it is a preprint. The 96.03\%
post-QC completeness figure is the number a network operator should take away.}
{10.21203/rs.3.rs-10441248/v1}

\paperentry{Chaiwong et al.\ 2026}
{PM$_{2.5}$ and systemic inflammatory biomarkers in elderly COPD versus healthy elderly}
{Int J Mol Sci}
{Pilot panel of 38 COPD patients (73.6\,$\pm$\,7.1\,y) and 20 age-matched controls,
sampled in a high- (79.1\,$\pm$\,23.3\,$\mu$g/m$^3$) and a low-pollution
(13.9\,$\pm$\,3.6\,$\mu$g/m$^3$) Chiang Mai season, measuring hsCRP, IL-6, IL-8 and
TNF-$\alpha$. IL-8 was higher in COPD across both periods, but \textbf{no significant
within-group increase accompanied the 5.7-fold seasonal exposure contrast} --- a null the
authors attribute to inhaled corticosteroids (65.8\% of COPD) and statins (42\% COPD,
65\% controls). That explanation is untestable at $n=58$; equally consistent are ceiling
effects on chronic inflammation, ambient-monitor exposure misclassification, and simple
lack of power. Reported as a pilot, and should be cited as one. Relevance: a seasonal
contrast this large failing to move systemic biomarkers argues for personal rather than
municipal exposure metrics in susceptible-population panels.}
{10.3390/ijms27146283}

\band{Mechanistic toxicology\hfill 2 records}{Violet}

\paperentry{Fakfum et al.\ 2026}
{Seasonal PM$_{2.5}$ exposure and plasma metabolome changes related to metabolic syndrome}
{Toxics}
{Prospective within-subject comparison of plasma $^1$H-NMR metabolomes in 25 healthy
adults in Samoeng District, Chiang Mai, between haze and non-haze seasons.
Twenty-six metabolites separated the seasons by PLS-DA (VIP $>$ 1.5): eleven up
(maleylacetoacetic acid, deoxyribose 5-phosphate, betaine, 3-hydroxyanthranilic acid,
1-methyladenosine) and fifteen down (deoxyguanosine, D-arabitol, glycerophosphocholine,
ophthalmic acid, oxaloacetic acid), implicating inflammation, ROS generation, NO
inhibition and disrupted energy metabolism across 56 pathways. The design cannot separate
PM$_{2.5}$ from everything else that changes with the Thai dry season --- temperature,
humidity, diet, activity, and co-emitted gases --- and PLS-DA at $n=25$ over hundreds of
features is prone to overfitting without the permutation testing the abstract does not
report. Hypothesis-generating.}
{10.3390/toxics14070544}

\paperentry{Lu et al.\ 2026}
{Toxicity of PM$_{2.5}$ from incomplete solid-fuel burning in \textit{Caenorhabditis elegans}}
{Toxics}
{Synchronised L4 \textit{C.\ elegans} exposed to aqueous suspensions of PM$_{2.5}$
collected from incomplete combustion of six solid fuels (rice, wheat, peanut and rapeseed
straw; poplar and paulownia branch), scored for body length, fertilised-egg count,
lipofuscin accumulation and ROS, alongside OC/EC, water-soluble ion and PAH
characterisation. \textbf{Oilseed-crop straw (peanut, rapeseed) was the most toxic}, and
the effect tracked lipid-soluble PAH content via oxidative stress. The source-specificity
result is the useful part --- it says fuel type, not mass concentration, drives toxicity,
which is an argument for speciation-capable monitoring. The model is the limitation: an
aqueous suspension delivered to a nematode cuticle and gut has no relationship to
pulmonary deposition, and no mammalian or human-cell confirmation is offered.}
{10.3390/toxics14070597}

\band{Occupational \& indoor exposure\hfill 2 records}{Clay}

\paperentry{Komonnithiphong et al.\ 2026}
{PM$_{2.5}$-bound PAHs and carcinogenic risk in charcoal barbecue (Moo Kratha) restaurants}
{Toxics}
{Ten area samples taken at customer dining tables in ten Chiang Mai charcoal-barbecue
restaurants during peak evening hours in February 2026 (early dry season), 2\,L\,min$^{-1}$
for 180\,min onto quartz fibre filters, 16 US-EPA priority PAHs by GC-MS. Mean PM$_{2.5}$
\textbf{87.64\,$\mu$g\,m$^{-3}$}, total PAHs 132.74\,ng\,m$^{-3}$, carcinogenic fraction
52.3\%, TEQ$_{\mathrm{BaP}}$ 15.34\,ng\,m$^{-3}$. Inhalation ILCR for a frequent adult
diner \textbf{2.82\,$\times$\,10$^{-5}$} (child 2.12\,$\times$\,10$^{-5}$), rising to
6.16 and 4.62\,$\times$\,10$^{-5}$ at the worst site; even monthly diners exceed
10$^{-6}$. Two hard limitations: $n=10$ single visits with no repeat sampling gives no
within-site variance, and February in Chiang Mai means the regional haze background is
confounded with the cooking source --- no outdoor or upwind reference was reported.
A paired indoor/outdoor low-cost node would have resolved that for a fraction of the cost.}
{10.3390/toxics14070643}

\paperentry{Subway PSD study 2026}
{PM and CO$_2$ in Chinese subway stations with varying platform screen door forms}
{J Environ Sci (China)}
{Seven stations instrumented across platform, hall and outdoor locations, contrasting
full-height platform screen doors (PSD) with half-height platform bailout doors (HPBD).
Platform $>$ hall $>$ outdoor for PM; CO$_2$ tracks passenger count linearly. The piston
effect on train arrival raises \textbf{PM$_{2.5}$ by 20\% and PM$_{10}$ by 36\% under
HPBD}, appreciably more than under PSD, while CO$_2$ falls 1.8\% (PSD) and 4.6\% (HPBD).
Fine particles carry 38.22--62.57\% of total mass, highest behind PSD, and Fe dominates
the elemental profile across six apportioned sources. The design is a snapshot without
seasonal or diurnal repetition, and the door form is confounded with station age, depth
and ventilation design, so the PSD-versus-HPBD contrast is associational. Still the
cleanest demonstration in today's corpus that a metric-and-geometry-specific enclosed
microenvironment needs its own sensor siting logic --- a platform-edge node and a
concourse node are measuring different populations of particles.}
{10.1016/j.jes.2025.12.027}

\band{Sensing, forecasting \& exposure modelling\hfill 2 records}{Amber}

\paperentry{Wang et al.\ 2026}
{Joint monitoring and early warning of SARS-CoV-2 in outdoor PM$_{2.5}$ and wastewater}
{Viruses}
{Twelve months (2023) of paired outdoor PM$_{2.5}$ filter and wastewater sampling in
Fuzhou, RT-qPCR for SARS-CoV-2 RNA, associations with pollutants and meteorology by GAM
and DLNM. \textbf{45.9\% of PM$_{2.5}$ samples were positive}; viral load in PM$_{2.5}$
rose with CO, PM$_{2.5}$ mass and atmospheric pressure and fell with temperature and
sunshine duration, with the \textbf{strongest lead over outpatient COVID cases at
$\approx$\,2\,days} against lag 0 for wastewater. The direct interest for a
distributed-sensing group is that a routine gravimetric filter stream doubles as a
bioaerosol surveillance substrate with a two-day lead --- but only if filters are archived
and chain-of-custody preserved, which most low-cost optical nodes cannot do at all. The
weakness is that RNA detection on a filter conflates aerosol transmission with resuspended
or deposited material, and no infectivity assay was attempted; the ``early warning'' claim
rests on ecological time-series correlation with outpatient counts.}
{10.3390/v18070788}

\paperentry{Ainembabazi et al.\ 2026}
{Trends in atmospheric pollution research in Uganda, 1990--2025: a scoping review}
{Toxics}
{Structured scoping review across Google Scholar, Web of Science and PubMed. Output was
negligible before 2010, expanded after 2015 and surged 2020--2025. PM$_{2.5}$/PM$_{10}$
dominate throughout while NO$_2$, SO$_2$, CO and O$_3$ are underrepresented; studies
cluster in urban Kampala with almost no rural ambient monitoring; and the methodological
trajectory runs from proxy and gravimetric approaches to \textbf{low-cost sensors,
portable monitors and satellite-derived products}. The authors' own conclusion is the
operative one: short-term, spatially constrained campaigns predominate, leaving no
long-term or nationally representative record. This is the clearest statement in today's
corpus of the niche a sustained calibrated low-cost network fills --- and a reminder that
Google Scholar inclusion makes the sampling frame non-reproducible, so the counts should
be treated as indicative.}
{10.3390/toxics14070542}

\paperentry{Paj\k{a}k et al.\ 2026}
{Toxicological legacy of PAHs from an urban tire fire: soil contamination and cancer risk}
{Toxics}
{Source-to-receptor assessment of a large landfill tire fire coupling HYSPLIT dispersion
modelling against air-quality monitor data with soil sampling on the fireground and in
surrounding allotment gardens, screened for Zn as a tire tracer and 16 priority PAHs.
Emissions \textbf{PM$_{10}$ 1.34\,t and PAHs 17.7\,kg}; fireground soil
$\Sigma$PAH 148.9\,mg\,kg$^{-1}$; cumulative child ILCR \textbf{9.7\,$\times$\,10$^{-4}$},
an order of magnitude above the 10$^{-4}$ regulatory benchmark, with dermal contact rather
than inhalation the dominant pathway and elevated risk (10$^{-5}$--10$^{-4}$) persisting
at distal residential sites. Ecological risk quotients exceeded unity for PAHs everywhere
and for Zn/Cu near the fire. Limitation: no pre-fire soil baseline, so urban background
PAH --- itself substantial in a coal-heating region --- is attributed to the fire; and
HYSPLIT at this scale is validated against a sparse monitor network. The transferable
point is that acute source events deposit a multi-year soil legacy that airborne
monitoring alone will not detect.}
{10.3390/toxics14070543}

\paperentry{Chengdu supersite 2026}
{Gas-phase reactions accelerate nitrate formation in spring and early summer, Sichuan Basin}
{J Environ Sci (China)}
{Urban supersite speciation of PM$_{2.5}$ over April--June 2018--2021 in Chengdu,
targeting the under-studied warm-season nitrate problem. During complex PM$_{2.5}$/O$_3$
episodes nitrate rose \textbf{47.8\%, from 5.52 to 8.16\,$\mu$g\,m$^{-3}$}, with its
share of PM$_{2.5}$ climbing from 13.4\% to 22.2\%. The NO$_3^-$/PM$_{2.5}$ ratio scales
with aerosol liquid water, and rapid early-morning ALW growth drives hygroscopic growth
and heterogeneous production on top of the gas-phase channel. Directly consequential for
optical sensing: a nitrate- and water-dominated aerosol has a very different hygroscopic
growth factor and refractive index from the winter carbonaceous mixture most low-cost
sensor calibrations are fitted to, so a humidity correction trained in winter will bias
warm-season estimates. Observational and single-city, so the mechanism is inferred from
covariation rather than tested.}
{10.1016/j.jes.2026.02.002}

\band{Burden, policy \& mitigation\hfill 3 records}{Coral}

\paperentry{Hu et al.\ 2026}
{Older-adult life expectancy in China: Bayesian mortality estimation and spatially heterogeneous determinants}
{Front Public Health}
{Two-stage design on the 2015 China 1\% Population Sample Survey: Bayesian random-effects
smoothing of age-specific mortality for 60+ across 31 provinces, abridged life tables, then
a bake-off of MLR, spatial error, spatial lag, GWR and MGWR for socio-environmental
determinants. \textbf{PM$_{10}$ was not significant in any global model}; only GWR
suggested a more negative association in south-western China, and GWR outperformed MGWR
($R^2$ 0.789 vs 0.664), which the authors correctly attribute to bandwidth instability at
$n=31$. Read this as a negative result about method, not about PM: 31 aggregated provincial
units, one cross-section, provincial mean PM$_{10}$ as the exposure, and no lag structure
cannot recover a pollution--longevity gradient that individual-level cohorts detect
routinely. The paper is a clean illustration of how ecological aggregation destroys an
exposure contrast that a denser monitoring grid would preserve.}
{10.3389/fpubh.2026.1851714}

\paperentry{Zhou et al.\ 2026b}
{Sources of ultrafine particles during severe haze periods: a spatial Durbin analysis}
{Toxics}
{Spatial Durbin model over the 28 ``2+26'' Beijing-Tianjin-Hebei channel cities, relating
pollution-intensive industrial activity to haze, and attributing ultrafine particle
generation to desulfurisation, denitrification and dust-removal operations --- specifically
lax wet-FGD management, ratcheting NO$_x$ standards, and electrostatic precipitators that
are ineffective in the ultrafine size range. Direct effect 0.028, total spillover-inclusive
effect 0.151. \textbf{No ultrafine particles were measured anywhere in this study.}
The UFP attribution is inferred from city-level industrial statistics and a haze proxy, so
the central claim --- that abatement equipment is itself a UFP source --- is an untested
hypothesis dressed as a finding, however physically plausible the ESP argument is. Logged
tier C for the hypothesis, not the evidence; it is exactly the claim a size-resolved
mobility-sizer campaign near a controlled stack could settle.}
{10.3390/toxics14070588}

\paperentry{Sun et al.\ 2026}
{Clarifying causal and translational inference in studies of ambient PM$_{2.5}$ and lung cancer}
{Br J Cancer}
{Correspondence, not primary research, arguing that the PM$_{2.5}$--lung-cancer literature
conflates aetiological causal claims with translational or prognostic ones, and that the
distinction changes what an estimate licenses. No new data. Worth keeping in the register
because the same conflation runs through exposure-assignment practice: a satellite- or
LUR-derived annual mean supports population-attributable-fraction arithmetic but not the
individual-level counterfactual that translational framing implies, and adding a denser
sensor network narrows the measurement error without touching the identification problem.}
{10.1038/s41416-026-03555-2}

\band{Other clinical endpoints\hfill 3 records}{Slate}

\paperentry{He et al.\ 2026}
{PM$_{2.5}$-bound chlorinated paraffins and sleep disorders in youth, with body weight as mediator}
{Toxics}
{Largest study in today's corpus: 122{,}965 valid questionnaires from school-aged children
across the Pearl River Delta, with PM$_{2.5}$-bound chlorinated paraffin exposure analysed
by generalised linear mixed models, restricted cubic splines, WQS regression and quantile
g-computation, plus mediation analysis. Per-IQR $\Sigma$CP exposure raised sleep-disorder
odds with \textbf{ORs reported as a 1.08--3.20 range}; short-chain CPs were the dominant
contributor; body weight mediated 3.40--35.67\%; associations were stronger in girls and
in children under 12. Two serious problems. The exposure is almost certainly modelled or
area-assigned rather than personally measured --- CP speciation on PM$_{2.5}$ is a
specialist analysis not deployable at $n>10^5$ --- so the exposure contrast is spatial and
confounded with everything else that varies across the PRD. And an OR range spanning
1.08--3.20 without a stated model set is uninterpretable: it could be lag structure,
subgroup, or CP congener. Large $n$ does not fix either.}
{10.3390/toxics14070607}

\paperentry{Ion-Andrei et al.\ 2026}
{Chronic air pollution exposure and the immunomolecular profile of lung cancer}
{Biomedicines}
{Regional retrospective series in Constan\c{t}a County, Romania, linking $\approx$15-year
historical monitor-based pollution exposure to tumour histology, PD-L1 expression and EGFR
status, with ALK/KRAS described only. \textbf{Chronic NO$_2$ was associated with higher
PD-L1 expression (OR 1.035, $p=0.017$)}, with PM$_{10}$, NO and NO$_x$ trending in the
same direction and an exploratory CO--EGFR signal. The idea --- that ambient exposure shapes
tumour immunophenotype and therefore immunotherapy eligibility --- is genuinely interesting
and underexplored. The execution does not support it: no CI or sample size in the
abstract, no adjustment for smoking described, an exposure assigned from county monitors
over 15 years with no residential-history weighting, and multiple correlated pollutants
tested without correction. Hypothesis-generating at best, and the authors concede
causality cannot be established.}
{10.3390/biomedicines14071615}

\paperentry{Do\u{g}an Karata\c{s} et al.\ 2026}
{Air pollution and sudden sensorineural hearing loss: nine-year retrospective study}
{J Clin Med}
{223 SSHL patients (mean 52.7\,$\pm$\,7.4\,y) at a single ENT clinic in Sivas, T\"urkiye,
2016--2025, with monthly mean PM$_{10}$, PM$_{2.5}$, SO$_2$, NO$_x$, NO, CO and NO$_2$ from
ministry monitoring stations and a ``cumulative risk'' computed by a public-health
specialist. Post-treatment audiometric change was significant at all frequencies
(greatest at 2\,kHz) and no significant seasonal variation in case occurrence was found.
The design cannot answer the question it poses: a case series with no control period and
no non-case comparison cannot estimate an exposure--occurrence association, and pairing
223 events with monthly city-wide pollutant means guarantees exposure misclassification
far exceeding any plausible effect. The audiometric outcomes are a treatment-response
analysis with no pollution content. Included for completeness of the register.}
{10.3390/jcm15145687}

\clearpage

% ==================================================================== TABLE
\section*{\sffamily\bfseries\color{Deep}Corpus register}
\vspace{-1.5mm}{\color{Amber}\rule{\linewidth}{1.1pt}}\vspace{2.5mm}

{\fontsize{6.7}{8.4}\selectfont
\setlength{\tabcolsep}{2.4pt}
\begin{longtable}{@{}>{\RaggedRight}p{25mm}>{\RaggedRight}p{41mm}>{\RaggedRight}p{27mm}%
>{\RaggedRight}p{23mm}>{\RaggedRight}p{18mm}>{\RaggedRight}p{15mm}@{}}
\toprule
\textsf{\textbf{First author}} & \textsf{\textbf{Focus}} & \textsf{\textbf{Journal}} &
\textsf{\textbf{Design}} & \textsf{\textbf{Metric}} & \textsf{\textbf{Region}}\\
\midrule
\endfirsthead
\toprule
\textsf{\textbf{First author}} & \textsf{\textbf{Focus}} & \textsf{\textbf{Journal}} &
\textsf{\textbf{Design}} & \textsf{\textbf{Metric}} & \textsf{\textbf{Region}}\\
\midrule
\endhead
\multicolumn{6}{@{}l@{}}{\cellcolor{Mist}\textsf{\textbf{\textcolor{Deep}{Neuro / mental health}}}}\\[0.6mm]
\href{https://doi.org/10.1001/jamanetworkopen.2026.25015}{Dales} & Wildfire-sourced PM$_{2.5}$ and ED visits for migraine and other primary headache & \textit{JAMA Netw Open} & Case-crossover & PM$_{2.5}$ (wildfire) & Canada\\
\href{https://doi.org/10.1093/ageing/afag215}{Glover} & Umbrella meta-review of urban environmental exposures and cognitive health & \textit{Age Ageing} & Umbrella review / meta-analysis & PM$_{2.5}$ & Global\\
\href{https://doi.org/10.1016/j.nsa.2026.107019}{Goossens} & In utero air pollution exposure and pre-/postnatal brain development & \textit{Neurosci Appl} & Narrative review & PM$_{2.5}$ / PM$_{10}$ / PM$_{\mathrm{coarse}}$ & Global\\
\href{https://doi.org/10.1016/j.isci.2026.116617}{Ma} & PM$_{2.5}$-induced neurotoxicity and depression: exposure risks and mechanisms & \textit{iScience} & Narrative review & PM$_{2.5}$ + components & Global\\
\href{https://doi.org/10.3760/cma.j.cn112150-20250804-00752}{Pan} & Maternal PM exposure and offspring neurodevelopment: effects and mechanisms & \textit{Zhonghua Yu Fang Yi Xue Za Zhi} & Narrative review & PM (bulk) & China\\
\addlinespace[0.9mm]
\multicolumn{6}{@{}l@{}}{\cellcolor{Mist}\textsf{\textbf{\textcolor{Teal}{Cardiovascular \& metabolic}}}}\\[0.6mm]
\href{https://doi.org/10.3760/cma.j.cn112338-20251120-00843}{Huang} & PM$_{2.5}$ and dyslipidaemia in rural Chinese children; dietary effect modification & \textit{Zhonghua Liu Xing Bing Xue Za Zhi} & Cross-sectional (multistage) & PM$_{2.5}$ & China\\
\href{https://doi.org/10.32394/pe/218966}{Kazmierska} & Air pollution and cardiovascular disease: review of current data & \textit{Przegl Epidemiol} & Narrative review & PM$_{2.5}$ / PM$_{10}$ & Global\\
\href{https://doi.org/10.1007/s00125-026-06814-2}{Xi} & Air pollution, built environment and temperature vs. prediabetes phenotypes & \textit{Diabetologia} & Prospective cohort & PM$_{2.5}$ & Germany\\
\href{https://doi.org/10.1161/JAHA.125.046917}{Yang} & Short-term PM$_{2.5}$ and hospitalisation/ED use in atrial fibrillation & \textit{J Am Heart Assoc} & Case-crossover & PM$_{2.5}$ (wildfire-influenced) & USA\\
\addlinespace[0.9mm]
\multicolumn{6}{@{}l@{}}{\cellcolor{Mist}\textsf{\textbf{\textcolor{Sky}{Reproductive \& developmental}}}}\\[0.6mm]
\href{https://doi.org/10.1016/j.jes.2025.10.043}{Chen} & Prenatal air pollution, fetal growth trajectory phenotypes and child neurodevelopment & \textit{J Environ Sci (China)} & Prospective cohort & PM$_{2.5}$ / PM$_{10}$ & China\\
\href{https://doi.org/10.1016/j.ecoenv.2026.120566}{Huang} & Joint prenatal bisphenol + air-pollutant mixture exposure and preschool IQ & \textit{Ecotoxicol Environ Saf} & Prospective cohort (mixtures) & PM$_{2.5}$ / PM$_{10}$ & China\\
\href{https://doi.org/10.1007/s40572-026-00552-8}{Peterson} & Ambient air pollution, PAHs and pubertal development: critical review & \textit{Curr Environ Health Rep} & Critical review & PM$_{2.5}$ components & Global\\
\href{https://doi.org/10.1016/j.envres.2026.125313}{Xie} & Long-term PM$_{2.5}$ and timing of pubertal development in boys and girls & \textit{Environ Res} & Prospective cohort & PM$_{2.5}$ & China\\
\addlinespace[0.9mm]
\multicolumn{6}{@{}l@{}}{\cellcolor{Mist}\textsf{\textbf{\textcolor{Sage}{Respiratory \& allergic}}}}\\[0.6mm]
\href{https://doi.org/10.1007/s00484-026-03289-z}{Bayar Muluk} & Climate change and otolaryngology: airway disease impacts and emission reduction & \textit{Int J Biometeorol} & Narrative review & PM$_{2.5}$ & Global\\
\href{https://doi.org/10.4103/mgr.MEDGASRES-D-25-00164}{Liao} & Short-term air pollution, respiratory symptoms and lung function in asthma & \textit{Med Gas Res} & Panel study & PM$_{2.5}$ / PM$_{10}$ & China\\
\href{https://doi.org/10.3390/epidemiologia7040097}{Sanchez de Prada} & Ambient PM and asthma-coded admissions in children 0-3 y, Spain & \textit{Epidemiologia} & Ecological panel & PM$_{2.5}$ / PM$_{10}$ & Spain\\
\href{https://doi.org/10.1016/j.jhazmat.2026.143076}{Yang} & Allergenic vs. toxic risk of airborne seed hairs compared with ambient PM$_{2.5}$ & \textit{J Hazard Mater} & Field measurement & PM$_{2.5}$ + bioaerosol & China\\
\addlinespace[0.9mm]
\multicolumn{6}{@{}l@{}}{\cellcolor{Mist}\textsf{\textbf{\textcolor{Violet}{Mechanistic toxicology}}}}\\[0.6mm]
\href{https://doi.org/10.1016/j.jes.2025.12.074}{Ding} & PM$_{2.5}$ promotes vascular calcification via miR-27b-5p/HES1 signalling & \textit{J Environ Sci (China)} & In vitro + organoid + murine & PM$_{2.5}$ & China\\
\href{https://doi.org/10.1007/s00360-026-01698-5}{Maraschi} & Osmoregulatory disruption and metal bioaccumulation from settleable atmospheric PM & \textit{J Comp Physiol B} & Controlled exposure (ecotox) & Settleable APM & Brazil\\
\href{https://doi.org/10.1016/j.jes.2025.11.018}{Wu} & Macrophage NLRP3 inflammasome in real-ambient PM-induced hepatic glucose dysregulation & \textit{J Environ Sci (China)} & Real-ambient murine + co-culture & Real-ambient PM & China\\
\addlinespace[0.9mm]
\multicolumn{6}{@{}l@{}}{\cellcolor{Mist}\textsf{\textbf{\textcolor{Clay}{Occupational \& indoor}}}}\\[0.6mm]
\href{https://doi.org/10.3389/fmedt.2026.1796828}{Chen} & Nested suction instrument to cut intraoperative surgical-smoke PM$_{2.5}$ & \textit{Front Med Technol} & Bench / device experiment & PM$_{2.5}$ (surgical smoke) & China\\
\href{https://doi.org/10.1093/annweh/wxag055}{Cole} & Determinants of respirable crystalline silica exposure in Australian tunnelling & \textit{Ann Work Expo Health} & Observational exposure study & Respirable dust / RCS & Australia\\
\href{https://doi.org/10.1016/j.jhazmat.2026.143081}{Ding} & Stage-resolved hair Zn-isotope dynamics as a firefighting combustion-exposure tracer & \textit{J Hazard Mater} & Repeated-measures biomarker & Combustion-derived PM & China\\
\href{https://doi.org/10.3390/jox16040119}{Ramsden} & Xenobiotic hazards in aircraft cabin air & \textit{J Xenobiot} & Review + risk economics & Ultrafine particles & Global\\
\addlinespace[0.9mm]
\multicolumn{6}{@{}l@{}}{\cellcolor{Mist}\textsf{\textbf{\textcolor{Amber}{Sensing, forecasting \& instrumentation}}}}\\[0.6mm]
\href{https://doi.org/10.1016/j.jenvman.2026.130515}{Shi} & Adaptive spatiotemporal graph transformer for multi-site multi-step PM$_{2.5}$ forecasting & \textit{J Environ Manage} & Deep-learning model & PM$_{2.5}$ & China / India\\
\href{https://doi.org/10.3390/bios16070380}{Tsuchiya} & Seasonal and indoor-outdoor characteristics of fluorescent aerosol particles in offices & \textit{Biosensors} & Real-time sensor campaign & Bioaerosol / coarse & Japan\\
\addlinespace[0.9mm]
\multicolumn{6}{@{}l@{}}{\cellcolor{Mist}\textsf{\textbf{\textcolor{Amber}{Exposure assessment \& modelling}}}}\\[0.6mm]
\href{https://doi.org/10.1029/2025GH001621}{Solorzano-Araque} & Multiscale exposure modelling: LOTOS-EUROS + mobility surveys + particle toxicity & \textit{GeoHealth} & CTM + mobility integration & PM$_{2.5}$ / PM$_{10}$ & Colombia\\
\addlinespace[0.9mm]
\multicolumn{6}{@{}l@{}}{\cellcolor{Mist}\textsf{\textbf{\textcolor{Coral}{Burden, policy \& mitigation}}}}\\[0.6mm]
\href{https://doi.org/10.1002/hsr2.72886}{Karamzad} & Ischaemic heart disease burden in MENA 1990-2021; PM attributable fraction 28.2\% & \textit{Health Sci Rep} & GBD secondary analysis & Ambient PM & MENA\\
\href{https://doi.org/10.1371/journal.pone.0347989}{Rees} & HEPA filtration to reduce respiratory infection in care homes & \textit{PLoS One} & RCT process evaluation & Filtered PM & UK\\
\href{https://doi.org/10.1016/j.jes.2026.01.009}{Shan} & Policy-driven emission reductions and post-pandemic rebound, high-resolution inventories & \textit{J Environ Sci (China)} & Emission inventory & PM emissions & China\\
\addlinespace[0.9mm]
\multicolumn{6}{@{}l@{}}{\cellcolor{Mist}\textsf{\textbf{\textcolor{Slate}{Other clinical endpoints}}}}\\[0.6mm]
\href{https://doi.org/10.3390/dj14070440}{Lupita} & PM$_{2.5}$, preventive-care capacity and caries/edentulism burden across the EU & \textit{Dent J} & Ecological panel & PM$_{2.5}$ & EU-27\\
\href{https://doi.org/10.1016/j.jdin.2026.06.008}{Park} & Long-term PM exposure and risk of androgenetic alopecia & \textit{JAAD Int} & Nationwide cohort & PM$_{2.5}$ / PM$_{10}$ & South Korea\\
\href{https://doi.org/10.3171/2026.3.PEDS2615}{Sun} & Environmental exposure and area deprivation vs. paediatric hydrocephalus outcomes & \textit{J Neurosurg Pediatr} & Registry cohort + geospatial & PM$_{2.5}$ (satellite-derived) & USA\\
\addlinespace[0.9mm]
\bottomrule
\end{longtable}}


\vspace{4mm}
\begin{keybox}
{\sffamily\bfseries\footnotesize\color{Deep}Provenance and method}\\[1.2mm]
{\fontsize{7.2}{8.8}\selectfont
\textbf{Entry window} 28 July 2026, continuous with the 24--27 July window of the previous
issue; no gap. \textbf{Sources queried:} PubMed E-utilities against
\texttt{[EDAT]} with two term sets (particle/health, 16 hits; sensing, calibration and
exposure-assessment terms crossed with particle terms, 1 hit, itself a subset of the
first); Europe PMC against \texttt{CREATION\_DATE} 26--28 July (10 hits, 5 already
published in the 27 July issue and removed by the dedup key); Crossref
\texttt{from-index-date} (rejected as a source --- Crossref re-indexes historical records
continuously, so the filter returned works from 2007 onward and is unusable for
entry-date windowing); OpenAlex \texttt{from\_created\_date}
(\textbf{now returns HTTP 403, ``Plan upgrade required'' --- this filter has moved behind a
paid tier and the OpenAlex leg of the pipeline is broken until an alternative is found});
arXiv (the \texttt{submittedDate} range query returned an empty body on this run;
the unfiltered relevance query surfaced nothing newer than 6 July). A supplementary web
sweep for low-cost-sensor calibration work confirmed the entry-window gap is real rather
than a query artefact.\\[1.2mm]
\textbf{Deduplication:} PMID $\rightarrow$ lowercased DOI $\rightarrow$ SHA-1 of
normalised title + first author + year, checked against \texttt{claude/state/seen.json}
before any summarisation. \textbf{Screening:} 21 candidates, 19 in scope. Two exclusions
logged to \texttt{claude/state/rejected.jsonl}: a deep-mine battery-fleet transferability
review (diesel PM as motivating context only, no measurement, exposure model or health
endpoint) and a water-transparency review (aquatic suspended particulate matter, not
airborne). \textbf{Labels} for subtopic, design, particle metric and geography are
single-label judgements from title, abstract and MeSH, not a controlled vocabulary.
\textbf{Effect estimates} are transcribed verbatim; exposure contrasts are heterogeneous
and estimates are not mutually comparable. One record is a \textbf{preprint} and is
labelled as such throughout. All DOI links resolve to the publisher of record; abstracts
remain the copyright of their respective publishers.\\[1.2mm]
\textbf{Correction, reissued 31 July 2026.} As first published, the Sun et al.\ entry
(PMID 42509466) carried \texttt{10.1016/j.lanwpc.2024.101106}, a DOI that PubMed's source
metadata had registered incorrectly and which resolved to an unrelated
\textit{Lancet Reg Health West Pac} article. The record has since been corrected upstream
and the link is now \texttt{10.1038/s41416-026-03555-2}
(\textit{Br J Cancer}). No other content in this issue has changed.\par}
\end{keybox}

\end{document}
